% Part: second-order-logic
% Chapter: sol-and-set-theory
% Section: comparing-sets

\documentclass[../../../include/open-logic-section]{subfiles}

\begin{document}

\olfileid{sol}{set}{cmp}

\olsection{مقارنة المجموعات}

\begin{prop}
!!{formula}~$\lforall[x][(X(x) \lif Y(x))]$ تعريف لعلاقة الاحتواء
الجزئي؛ فـ$\Sat{M}{\lforall[x][(X(x) \lif Y(x))]}[s]$ إذا وفقط
إذا~$s(X) \subseteq s(Y)$.
\end{prop}

\begin{prop}
!!{formula}~$\lforall[x][(X(x) \liff Y(x))]$ تعريف لهوية المجموعات؛
فـ$\Sat{M}{\lforall[x][(X(x) \liff Y(x))]}[s]$ إذا وفقط
إذا~$s(X) = s(Y)$.
\end{prop}

\begin{prop}
!!{formula}~$\lexists[x][X(x)]$ تعريف لعدم الخلو؛
فـ$\Sat{M}{\lexists[x][X(x)]}[s]$ إذا وفقط إذا~$s(X) \neq \emptyset$.
\end{prop}

قولنا إن المجموعة~$X$ لا تزيد حجمًا على المجموعة~$Y$، أي
$\cardle{X}{Y}$، معناه وجود دالة !!a{injective}~$f\colon X \to Y$.
ولما أمكن التعبير عن تباين الدالة وعن وقوع قيمها عند عناصر~$X$
في~$Y$، أمكن تعريف علاقة عدم الزيادة في الحجم بين مجموعات المجال
الجزئية.

\begin{prop}
الصيغة
\[
\lexists[u][(\lforall[x][(X(x) \lif Y(u(x)))] \land \lforall[x][\lforall[y][(\eq[u(x)][u(y)] \lif
      \eq[x][y])]])]
\]
تعريف لعلاقة عدم الزيادة في الحجم.
\end{prop}

ومجموعتان متساويتان حجمًا، أو «متكافئتان عدديًا»،
$\cardeq{X}{Y}$، إذا وفقط إذا وجدت بينهما دالة
!!a{bijective}~$f\colon X \to Y$.

\begin{prop}
الصيغة
\begin{multline*}
  \lexists[u][(\lforall[x][(X(x) \lif Y(u(x)))] \land {}\\
    \lforall[x][\lforall[y][((X(x) \land X(y) \land \eq[u(x)][u(y)]) \lif \eq[x][y])]]
  \land {} \\\lforall[y][(Y(y) \lif \lexists[x][(X(x)
      \land \eq[y][u(x)])])])]
\end{multline*}
تعريف لعلاقة التكافؤ العددي.
\end{prop}

وسنستعمل اختصارًا لهذه !!{formula}s الرموز
$X \subseteq Y$، و$X = Y$، و$X \neq \emptyset$، و$\cardle{X}{Y}$،
و$\cardeq{X}{Y}$، بهذا الترتيب.
وقد يوقع هذا الاستعمال في التباس؛ فالرموز نفسها نستعملها حين
نتكلم خارج اللغة الصورية على المجموعتين~$X$ و$Y$.
أما هنا فهي اختصارات لـ!!{formula}s من الرتبة الثانية،
ترد فيها متغيرات العلاقات الأحادية الموضع~$X$ و$Y$.

\begin{prop}
!!{sentence}~$\lforall[X][\lforall[Y][((\cardle{X}{Y} \land
    \cardle{Y}{X}) \lif \cardeq{X}{Y})]]$ صحيحة منطقيًا.
\end{prop}

\begin{proof}
إرضاء !!{sentence} في !!a{structure}~$\Struct{M}$
يعني أنه لكل مجموعتين جزئيتين~$X \subseteq \Domain{M}$
و$Y \subseteq \Domain{M}$ يحصل~$\cardeq{X}{Y}$ متى
حصل~$\cardle{X}{Y}$ و$\cardle{Y}{X}$.
وهذا ثابت لـ\emph{أي} مجموعتين~$X$ و$Y$ بمبرهنة شرودر--برنشتاين.
\end{proof}

\end{document}
